
\chapter{Die IBM-Regel}

\begin{ibmbox}
\textbf{I}nshallah, \textbf{B}ukra, \textbf{M}aalesh -- die drei Wörter,
die gemeinsam die gesamte ägyptische Zeitplanung erklären. Wenn ein Termin
,,inshallah bukra`` (so Gott will, morgen) stattfindet und es doch nicht
klappt, sagt man einfach ,,ma3lesh`` (macht nichts) und macht weiter.
Willkommen im System.
\end{ibmbox}

\begin{itemize}[leftmargin=*]
\item \arword{يلا}{yalla}{Los geht's! / Komm schon! / Beeil dich!}{%D9%8A%D9%84%D8%A7}
\item \arword{خلاص}{khaláS}{Fertig. Schluss. Genug. (das nützlichste Wort überhaupt)}{%D8%AE%D9%84%D8%A7%D8%B5}
\item \arword{معلش}{ma3lesh}{Macht nichts / Kein Problem / Tut mir leid}{%D9%85%D8%B9%D9%84%D8%B4}
\item \arword{إن شاء الله}{inshallah}{So Gott will (,,vielleicht, vielleicht auch nicht``)}{%D8%A5%D9%86%20%D8%B4%D8%A7%D8%A1%20%D8%A7%D9%84%D9%84%D9%87}
\item \arword{بكرة}{bükra}{Morgen (der Tag nach heute -- theoretisch)}{%D8%A8%D9%83%D8%B1%D8%A9}
\item \arword{حاضر}{HáaDir}{Wird gemacht! / Jawohl! (nützlich gegenüber Am Mahmoud)}{%D8%AD%D8%A7%D8%B6%D8%B1}
\end{itemize}

\begin{lachbox}
\textarabic{خلاص} (khal\'aS) ist das Schweizer Taschenmesser der
arabischen Sprache. Es bedeutet: ,,fertig``, ,,genug jetzt``, ,,okay,
abgemacht``, ,,lass uns aufhören zu streiten`` -- und manchmal einfach nur
,,...``. Wenn du nur ein Wort aus diesem Buch behältst, nimm dieses.
\end{lachbox}

\begin{bissobox}
Wenn Bisso genug vom Streicheln hat, sagt sie es nicht mit Worten -- sie
geht einfach. Julia, das ist \textarabic{خلاص} in Katzenform.
\end{bissobox}

\begin{tippbox}
\textarabic{بكرة} (bukra, ,,morgen``) ist mit Vorsicht zu genießen -- es
bedeutet nicht zwingend ,,in 24 Stunden``. Es bedeutet eher ,,nicht heute,
aber irgendwann bestimmt``. Plane entsprechend.
\end{tippbox}

\begin{checkpointbox}
Baue einen Satz mit allen drei IBM-Wörtern: ,,Kommst du morgen?`` --
,,Inshallah, bukra!`` -- und wenn es nicht klappt: ,,Ma3lesh!`` Übe das
laut, am besten mit einem Augenzwinkern.
\end{checkpointbox}
